% GENERATED FILE. Edit canonical lessons, then run npm run generate:books.
% canonical-source-hash: fnv1a64:4512788f890a7ca1
% canonical-lessons: PT-C259-falava, PT-C259-tinha-havia, PT-C259-fazia-frio, PT-C259-dormia-dormi, PT-C259-estava-a, PT-C259-revisao-quando

\chapter{The Scene --- Era, Tinha, Havia}
\label{ch:pt-the-scene}
\hlchaptermodality{259}

\begin{chapteropening}
\textbf{By the end of this chapter:} \emph{I can paint the scene of a past story with the pretérito imperfeito, and drop the events into it with the pretérito perfeito.}

Complete PT-C259-revisao-quando: a scene in the imperfeito and an event that cuts into it, joined with quando.
\end{chapteropening}

\section[falava, comia, partia]{falava — the past that is still running}

\label{lesson:PT-C259-falava}



\begin{quote}
Say \textquotedblleft{}I spoke\textquotedblright{} the way chapter fifteen built it. (\emph{Falei.})

That past has an edge: it happened, and it finished. Portuguese has a second past with no edge at all, and every story you tell from here on needs both of them.
\end{quote}

\begin{grammarlens}[title={Grammar Lens: two endings, three families}]
Take the infinitive and drop its last two letters. An \emph{-ar} verb adds \textbf{-ava}; an \emph{-er} or an \emph{-ir} verb adds \textbf{-ia}:

\begin{quote}
fal-ar, \textbf{falava} — I was speaking, I used to speak com-er, \textbf{comia} — I was eating, I used to eat part-ir, \textbf{partia} — I was leaving
\end{quote}

The other persons keep that ending and change only what comes after it: \textbf{falava}, \textbf{falavas}, \textbf{falava}, \textbf{falávamos}, \textbf{falavam} — I, you, he or she, we, they. And for the other two families: \textbf{comia}, \textbf{comias}, \textbf{comia}, \textbf{comíamos}, \textbf{comiam}.

Two things to hear. The \emph{nós} form carries a written accent, \emph{falávamos}, \emph{comíamos}, because the stress sits further back than usual. And the forms for \textquotedblleft{}I\textquotedblright{} and for \textquotedblleft{}he or she\textquotedblright{} are the same word: \emph{eu falava}, \emph{ele falava}. When it matters, keep the \emph{eu}.

This is the \textbf{pretérito imperfeito}, and almost every verb follows this one recipe. The few that refuse begin in the next lesson.

What it means is \textquotedblleft{}was ‑ing\textquotedblright{}, or \textquotedblleft{}used to\textquotedblright{}: \emph{morava em Lisboa}, I lived in Lisbon, back then, with no start and no finish in view. That is exactly what a story needs for its background.
\end{grammarlens}

\begin{cousinweb}
Latin already had this past, and it marked it with a \textbf{b}: Latin \emph{amābat}, \textquotedblleft{}he was loving\textquotedblright{}, and \emph{habēbat}, \textquotedblleft{}he was having\textquotedblright{}. In an \emph{-ar} verb that \emph{b} softened to a \textbf{v}, so \emph{falava} goes back to a spoken Latin past of \emph{fābulārī}, \textquotedblleft{}to talk\textquotedblright{}, the verb that gave \textbf{falar} (EtymDB 2.1). After \emph{e} and \emph{i} the \emph{b} wore away entirely, which is why \emph{comia} and \emph{partia} end in a bare \emph{-ia}.

Neither EtymDB nor Lewis \& Short records endings this late, so that last step is the usual account in histories of Portuguese, given here as one.
\end{cousinweb}

\subsection*{Your turn}
Say these aloud:

\begin{itemize}
\raggedright
  \item \emph{falava, comia, partia} — one ending for \emph{-ar}, one for the rest
  \item \emph{falava, falavas, falava, falávamos, falavam}
  \item \emph{trabalhava} — chapter five's verb, still running in the past
  \item the edge, then no edge — \emph{falei}, then \emph{falava}
\end{itemize}

\begin{tcolorbox}[breakable,colback=teal!4,colframe=teal!35!black,title={Before you move on}]
How do you build \textquotedblleft{}I was eating\textquotedblright{}? (Drop the \emph{-er} of \emph{comer} and add \emph{-ia}: \emph{\textbf{comia}}.) Which Latin letter became the \emph{-v-} of \emph{falava}? (The \emph{\textbf{b}} of \emph{amābat}.) Which past has an edge, and which has none? (\emph{\textbf{Falei}} has one; \emph{\textbf{falava}} has none.)
\end{tcolorbox}
\section[tinha, havia]{tinha, havia — the words a scene is made of}

\label{lesson:PT-C259-tinha-havia}



\begin{quote}
Build \textquotedblleft{}I was speaking\textquotedblright{} from \emph{falar}. (\emph{Falava.}) Say \textquotedblleft{}I am twenty\textquotedblright{} the way chapter fourteen taught it. (\emph{Tenho vinte anos.}) Then \textquotedblleft{}there is bread\textquotedblright{}. (\emph{Há pão.})

Hold on to all three. This lesson moves them into the past.
\end{quote}

\subsection*{You'll want to know: era and tinha}
Only four verbs refuse last lesson's recipe, and two of them are the ones you use most. The first you met in chapter sixteen: \textbf{ser} gives \textbf{era}.

\begin{quote}
\textbf{era}, \textbf{eras}, \textbf{era}, \textbf{éramos}, \textbf{eram} \emph{A estação era pequena.} — The station was small.
\end{quote}

The second is \textbf{ter}, and it gives \textbf{tinha}:

\begin{quote}
\textbf{tinha}, \textbf{tinhas}, \textbf{tinha}, \textbf{tínhamos}, \textbf{tinham} \emph{Tinha vinte anos.} — I was twenty. \emph{Tínhamos um cão.} — We had a dog.
\end{quote}

Age is the clearest case. \emph{Tenho vinte anos} is now; \emph{tinha vinte anos} is the age you were, all through some stretch of the past. The other two refusers are \textbf{vir}, which gives \emph{vinha}, and \textbf{pôr}, which gives \emph{punha}; you will meet them when you need them.

\subsection*{You'll want to know: havia}
You have \textbf{há}, \textquotedblleft{}there is\textquotedblright{} and \textquotedblleft{}there are\textquotedblright{}. Its past is \textbf{havia}, and it follows the recipe after all, since it comes from \emph{haver}:

\begin{quote}
\emph{Havia muitas pessoas.} — There were a lot of people. \emph{Não havia pão.} — There was no bread.
\end{quote}

Like \emph{há}, it does not change for a plural: one \emph{havia} for one thing or for a hundred. That is how a European Portuguese scene opens. In everyday Brazilian speech you will often hear \emph{tinha} doing this job instead: \emph{tinha muitas pessoas}.

\begin{cousinweb}
\textbf{ter} is inherited from Latin \emph{tenēre}, \textquotedblleft{}to hold\textquotedblright{} (EtymDB 2.1; Lewis \& Short, under \emph{teneo}), and \textbf{tinha} goes back to its past \emph{tenēbat}, \textquotedblleft{}he was holding\textquotedblright{}. \textbf{havia} goes back to \emph{habēbat}, \textquotedblleft{}he was having\textquotedblright{}, the past of \emph{habēre}, which gave \textbf{haver} (EtymDB 2.1). Its first \emph{b} softened to \emph{v}, as in \emph{falava}, and the second wore away, as in \emph{comia}.

\textbf{era} is Latin \emph{erat}, \textquotedblleft{}he was\textquotedblright{} (Lewis \& Short, under \emph{sum}). Latin's own past of \textquotedblleft{}to be\textquotedblright{} had no \emph{b}, which is why \emph{era} looks nothing like \emph{falava}. How the \emph{-nh-} of \emph{tinha} arose is not something either source traces, so it is left open here.
\end{cousinweb}

\subsection*{Your turn}
Say these aloud:

\begin{itemize}
\raggedright
  \item \emph{era, tinha, havia}
  \item \emph{Tinha vinte anos.} — then \emph{Tínhamos um cão.}
  \item \emph{Havia muitas pessoas.}
  \item from now to then — \emph{há} becomes \emph{havia}, \emph{tenho} becomes \emph{tinha}
\end{itemize}

\begin{tcolorbox}[breakable,colback=teal!4,colframe=teal!35!black,title={Before you move on}]
Which Latin verb is inside \emph{tinha}? (\emph{\textbf{Tenēre}}, \textquotedblleft{}to hold\textquotedblright{}.) Why does \emph{era} not end in \emph{-va}? (Latin's own form had no \emph{\textbf{b}}.) Say \textquotedblleft{}there were a lot of people\textquotedblright{}. (\emph{\textbf{Havia muitas pessoas.}}) And does \emph{havia} change for a plural? (\textbf{No.})
\end{tcolorbox}
\section[fazia frio, fazia calor]{fazia frio — the weather, made in the past}

\label{lesson:PT-C259-fazia-frio}



\begin{quote}
Say \textquotedblleft{}there were a lot of people\textquotedblright{} the way last lesson built it. (\emph{Havia muitas pessoas.}) Then \textquotedblleft{}I was twenty\textquotedblright{}. (\emph{Tinha vinte anos.})
\end{quote}

\subsection*{You'll want to know: fazia frio, fazia calor}
English says the weather \textbf{was} cold. Portuguese says it \textbf{made} cold:

\begin{quote}
\textbf{Fazia frio.} — It was cold. \textbf{Fazia calor.} — It was hot.
\end{quote}

\emph{Fazia} is the imperfeito of \textbf{fazer}, \textquotedblleft{}to make, to do\textquotedblright{}, on the plain \emph{-er} recipe: \emph{faz-er}, \emph{fazia}. \textbf{frio} is \textquotedblleft{}cold\textquotedblright{} and \textbf{calor} is \textquotedblleft{}heat\textquotedblright{}, so the second one says, word for word, \textquotedblleft{}it was making heat\textquotedblright{}. Neither needs a subject: nobody is doing the making.

Rain has a verb of its own. You have \textbf{a chuva}, \textquotedblleft{}the rain\textquotedblright{}, and its verb \textbf{chover} gives \textbf{chovia}, \textquotedblleft{}it was raining\textquotedblright{}, on the same recipe.

Keep \textbf{estava} for things, with the adjectives you already own. \emph{O café estava frio} — the coffee was cold — is about the coffee; \emph{estava quente}, it was hot, about a cup or a room. \emph{Fazia frio} is about the whole day. Put three of these together and you have a morning before anything has happened in it:

\begin{quote}
\emph{Fazia frio e chovia. Não havia sol.}
\end{quote}

\begin{cousinweb}
\textbf{fazer} is inherited from Latin \emph{facere}, \textquotedblleft{}to make, to do\textquotedblright{} (EtymDB 2.1), and \textbf{fazia} goes back to its past \emph{faciēbat}, \textquotedblleft{}he was making\textquotedblright{} (Lewis \& Short, under \emph{facio}). That \emph{facere} is the root of English \textbf{fact} and \textbf{factory}.

\textbf{frio} is inherited from Latin \emph{frīgidus}, \textquotedblleft{}cold\textquotedblright{} (EtymDB 2.1), worn short in the mouth, and English \textbf{frigid} borrowed the same word whole. \textbf{calor} is Latin \emph{calor}, \textquotedblleft{}heat\textquotedblright{}, kept almost letter for letter (EtymDB 2.1; Lewis \& Short). \textbf{chover} is inherited from spoken Latin \emph{plovere} (EtymDB 2.1), a reshaping of classical \emph{pluere}, \textquotedblleft{}to rain\textquotedblright{} (Lewis \& Short, under \emph{pluo}).
\end{cousinweb}

\subsection*{Your turn}
Say these aloud:

\begin{itemize}
\raggedright
  \item \emph{Fazia frio.} — then \emph{Fazia calor.}
  \item \emph{chovia} — \emph{chover} on the plain recipe
  \item a thing, then a day — \emph{O café estava frio.} then \emph{Fazia frio.}
  \item a whole morning — \emph{Fazia frio e chovia. Não havia sol.}
\end{itemize}

\begin{tcolorbox}[breakable,colback=teal!4,colframe=teal!35!black,title={Before you move on}]
Say \textquotedblleft{}it was hot\textquotedblright{}. (\emph{\textbf{Fazia calor.}}) Which do you use for the coffee, \emph{estava} or \emph{fazia}? (\emph{\textbf{Estava}} — it is about a thing.) Which Latin verb is inside \emph{fazia}, and what English word shares it? (\emph{\textbf{Facere}}, \textquotedblleft{}to make\textquotedblright{} — \textbf{fact}, \textbf{factory}.)
\end{tcolorbox}
\section[dormia / dormi]{dormia / dormi — the same sleep, told two ways}

\label{lesson:PT-C259-dormia-dormi}



\begin{quote}
Paint a cold morning in two short sentences. (\emph{Fazia frio. Não havia sol.})

Nothing has happened in it yet. This lesson is about what happens.
\end{quote}

\begin{grammarlens}[title={Grammar Lens: the edge for -er and -ir verbs}]
Chapter fifteen gave you the past with an edge for \emph{-ar} verbs: \emph{falei}, \emph{falou}. The \emph{-er} and \emph{-ir} verbs make it with an \textbf{-i}:

\begin{quote}
\textbf{dormi} — I slept · \textbf{dormiu} — he or she slept \textbf{comi} — I ate · \textbf{comeu} — he or she ate
\end{quote}

The other persons follow: \emph{dormiste}, \emph{dormimos}, \emph{dormiram}. Listen to the stress, which falls on that last \emph{-i}: \emph{dor-MI}. And set it beside the imperfeito: \textbf{dormi} has an edge, \textbf{dormia} has none. One letter apart, and they answer different questions.
\end{grammarlens}

\begin{grammarlens}[title={Grammar Lens: the scene and the events}]
Take one fact, that you slept, and let the tense decide what kind of fact it is:

\begin{quote}
\textbf{Dormia.} — I was sleeping. \emph{(no edge: what was going on)} \textbf{Dormi.} — I slept. \emph{(an edge: a thing that happened and ended)}
\end{quote}

Neither is more correct. A told story needs both, in a division of labour:

\begin{quote}
The \textbf{imperfeito} paints the \textbf{scene}: what was going on, what things were like. The \textbf{perfeito} drops the \textbf{events} into it: what happened.
\end{quote}

\emph{Fazia frio. Eu dormia. A minha mãe telefonou.} — It was cold. I was asleep. My mother phoned.

English has no separate ending for this. It leans on \emph{was ‑ing} and \emph{used to} and lets the context carry the rest. Portuguese settles it in the ending.
\end{grammarlens}

\begin{cousinweb}
\textbf{imperfeito} is inherited from Latin \emph{imperfectus}, \textquotedblleft{}unfinished\textquotedblright{} (EtymDB 2.1): \emph{in-}, \textquotedblleft{}not\textquotedblright{}, plus \emph{perfectus}, \textquotedblleft{}carried all the way through\textquotedblright{}, which is itself \emph{per-}, \textquotedblleft{}through\textquotedblright{}, plus \emph{facere}, \textquotedblleft{}to make\textquotedblright{} (Lewis \& Short). That \emph{facere} is the one inside last lesson's \emph{fazia}.

So the two names describe the two jobs. The \textbf{perfeito} is the past carried all the way through, to its edge. The \textbf{imperfeito} is the \textbf{not-carried-through} past: unfinished in time, with no edge on it.
\end{cousinweb}

\subsection*{Your turn}
Say these aloud:

\begin{itemize}
\raggedright
  \item the pair — \emph{dormia}, then \emph{dormi}
  \item \emph{comi, comeu} — then \emph{dormi, dormiu}
  \item a scene, then an event — \emph{Chovia. A minha mãe telefonou.}
  \item \emph{Havia pão, por isso comi.}
\end{itemize}

\begin{tcolorbox}[breakable,colback=teal!4,colframe=teal!35!black,title={Before you move on}]
In a story, which past paints and which drops the events in? (The \emph{\textbf{imperfeito}} paints; the \emph{\textbf{perfeito}} drops them in.) Say \textquotedblleft{}I ate\textquotedblright{}. (\emph{\textbf{Comi.}}) What does \emph{imperfectus} say? (\textbf{Not carried through.})
\end{tcolorbox}
\section[estava a dormir quando…]{estava a dormir quando… — the scene, and what cut into it}

\label{lesson:PT-C259-estava-a}



\begin{quote}
Say \textquotedblleft{}I was sleeping\textquotedblright{}, then \textquotedblleft{}I slept\textquotedblright{}. (\emph{Dormia}, then \emph{dormi}.)

There is a second way to say the first one, and in Portugal it is the way you will hear most.
\end{quote}

\begin{grammarlens}[title={Grammar Lens: estava a, and an infinitive}]
Take \textbf{estar}, put it in the imperfeito, add \textbf{a} and an infinitive:

\begin{quote}
\textbf{Estava a dormir.} — I was sleeping. \textbf{Estava a ler.} — I was reading. \textbf{Estávamos a falar.} — We were talking.
\end{quote}

\emph{Estar} follows the plain \emph{-ar} recipe here: \emph{estava}, \emph{estavas}, \emph{estava}, \emph{estávamos}, \emph{estavam}. Only the \emph{estar} changes; the infinitive after \emph{a} stays as it is.

\emph{Estava a ler} and \emph{lia} both paint the scene. \emph{Estava a} holds it still at one moment, a reading caught in progress; \emph{lia} can also mean \textquotedblleft{}I used to read\textquotedblright{}. The present works the same way: \emph{estou a ler}, I am reading right now.

This \emph{a} plus infinitive is the Portuguese norm. In Brazil, and in parts of southern Portugal, the same idea takes a form ending in \emph{-ndo}: \emph{estava dormindo}. Both are understood everywhere.
\end{grammarlens}

\begin{grammarlens}[title={Grammar Lens: quando, where the event comes in}]
\textbf{quando} means \textquotedblleft{}when\textquotedblright{}, and it puts a scene and an event into one sentence:

\begin{quote}
\emph{Estava a dormir \textbf{quando} a minha mãe telefonou.} — I was sleeping when my mother phoned. \emph{Estava a ler \textbf{quando} o telefone tocou.} — I was reading when the phone rang.
\end{quote}

The \emph{estava a} half is the scene, still running. The \emph{quando} half is the event that cut into it, in the perfeito: \emph{telefonou}, \emph{tocou}. You can turn the sentence round, and the tenses keep their jobs: \emph{Quando o telefone tocou, eu estava a ler.}
\end{grammarlens}

\begin{cousinweb}
\textbf{estar} is inherited from Latin \emph{stāre}, \textquotedblleft{}to stand\textquotedblright{} (EtymDB 2.1), so \emph{estava a ler} is, word for word, \textquotedblleft{}I was standing at reading\textquotedblright{}: caught in the middle of it. \textbf{quando} is inherited from Latin \emph{quandō}, \textquotedblleft{}when\textquotedblright{} (EtymDB 2.1; Lewis \& Short), almost unchanged.
\end{cousinweb}

\subsection*{Your turn}
Say these aloud:

\begin{itemize}
\raggedright
  \item \emph{Estava a dormir.} — then \emph{Estávamos a falar.}
  \item \emph{Estava a ler quando o telefone tocou.}
  \item the same, turned round — \emph{Quando o telefone tocou, eu estava a ler.}
  \item a scene of your own — \emph{Fazia frio e eu estava a …} — then one thing that happened, with \emph{quando}
\end{itemize}

\begin{tcolorbox}[breakable,colback=teal!4,colframe=teal!35!black,title={Before you move on}]
What goes between \emph{estava} and the infinitive? (\emph{\textbf{A.}}) In \emph{estava a ler quando o telefone tocou}, which half is the scene? (\emph{\textbf{Estava a ler.}}) Which past does the \emph{quando} half take here? (The \textbf{perfeito} — \emph{tocou}, the event.) And what does \emph{estar} mean, at its root? (\textbf{To stand.})
\end{tcolorbox}
\section[(review)]{Review — the scene, and what cut into it}

\label{lesson:PT-C259-revisao-quando}



\begin{quote}
Build \textquotedblleft{}I was reading\textquotedblright{} from \emph{ler}. (\emph{Lia}, or \emph{estava a ler}.) Then say \textquotedblleft{}I had a dog\textquotedblright{}. (\emph{Tinha um cão.})
\end{quote}

\begin{grammarlens}[title={What you've built}]
Five lessons, and every piece of a scene is in your hands. The recipe, \textbf{-ava} and \textbf{-ia}, on almost any verb: \emph{falava}, \emph{comia}, \emph{partia}. The four that refuse it, led by \emph{era} and \emph{tinha}. The scene-setter \emph{havia}. The weather, made rather than being: \emph{fazia frio}. The European \emph{estava a}, for a moment caught in progress. And the division of labour: the imperfeito paints, the perfeito drops the events in.

The joint between the two is \textbf{quando}.

\begin{quote}
\emph{Estava a ler quando a minha mãe telefonou.} \emph{Estava à espera do autocarro quando começou a chover.}
\end{quote}

The second one is worth hearing twice. The rain is an event here, because it \textbf{began}: \emph{começou}. Once it is falling, it is scene again: \emph{chovia}. And \emph{estar à espera}, \textquotedblleft{}to be waiting\textquotedblright{}, is the European way to say it, with \emph{espera} from your \emph{esperar}.
\end{grammarlens}

\subsection*{Your turn}
Say these aloud:

\begin{itemize}
\raggedright
  \item a scene in three — \emph{Fazia frio. Havia muitas pessoas. Eu tinha fome.}
  \item \emph{Estava a ler quando a minha mãe telefonou.}
  \item the rain as an event — \emph{começou a chover} — then as a scene — \emph{chovia}
  \item one of your own — \emph{Estava a …} plus one thing that happened, with \emph{quando}
\end{itemize}

\begin{tcolorbox}[breakable,colback=teal!4,colframe=teal!35!black,title={Before you move on}]
After \emph{quando}, in these sentences, which past? (The \textbf{perfeito} — the event.) Why is \emph{começou a chover} in the perfeito? (The rain \textbf{began}: that is an event.) What does \emph{imperfeito} mean, taken apart? (\textbf{Not carried through.}) And where did the \emph{-v-} of \emph{falava} come from? (The \emph{\textbf{b}} of Latin \emph{amābat}.)
\end{tcolorbox}
